\section{ハフマン符号に関する説明}
\subsection{歴史}
    \label{sec:Sec1Implementation}
ハフマン符号（Huffman Coding）は、1952年にDavid A. Huffmanによって考案された可逆データ圧縮アルゴリズムで、文字の出現頻度に基づいて可変長の符号を割り当て、データサイズを削減するために用いられた符号である
\subsection{基本原理}
頻出文字に短い符号を割り当て、稀な文字に長い符号を割り当てることで通常の全文字同じ長さでは実現できない保存手法。ただし、解凍時に復元可能である必要があり、たとえばA=0，B=1,C=10のような割当は不可(10はBAもしくはCであるため)。復元可能な最適な割り当てを作るためハフマン木(Huffman Tree)が考察された
\subsection{具体例}

入力テキスト:"ABRACADABRA"

文字頻度:

A: 5回 (45.5\%)

B: 2回 (18.2\%) 

R: 2回 (18.2\%)

C: 1回 (9.1\%)

D: 1回 (9.1\%)

ハフマン符号:

A: 0     (1ビット)

B: 10    (2ビット)

R: 11    (2ビット)

C: 110   (3ビット)

D: 111   (3ビット)

元のデータ: 11文字 × 8ビット = 88ビット

圧縮後: 5×1 + 2×2 + 2×2 + 1×3 + 1×3 = 20ビット

圧縮率: 77.3%


符号の割り当て方法:

頻度順に並ぶと、D1,C1,R2,B2,A5になり。その頻度の小さい順に新しいrootを作る。

そうするとD1,C1,R2,B2,A5 -- root(CD)2,R2,B2,A5 -- B2,root(CDR)4,A5 -- A5,root(CDRB)6--root(CDRBA)11になる

生成された木から逆に文字を探し、A5,root(CDRB)6--root(CDRBA)11の左部分を0右部分を1とすると、Aは0，root(CDRB)は1になる。B2,root(CDR)4--root(CDRB)6であるためB2はroot(CDRB)と0、つまり10、root(CDR)は11。

\subsection{可逆の原因}
Huffman Treeの構造上、全ての文字は葉の部分にあり、rootになっていないため他の文字の先頭になっていない
